\section{State-Resolved Element Properties}
\label{manual2:element-properties}

\paragraph{Setup.}
The property catalog distinguishes an element, an isotope in a nuclear state,
and an atom or ion in an electronic state. Its identity fields are atomic
number, isotope assignment, nuclear state, charge and electronic state. An
unspecified isotope remains unspecified. A neutral ground-state ionization
energy describes the transition to the singly charged continuum; the second
ionization energy has a singly charged initial state. Electron affinity also
requires the negative-ion endpoint. A nuclear mass and an atomic mass have
different electron inventories.

Each property record binds one entity, one quantity, its unit, source locator,
value qualification and uncertainty convention. The method of obtaining a
value, prior exposure to the target and the scope of a model claim are separate
fields. An evaluated value may incorporate a measurement or an extrapolation;
neither designation establishes a native AOD prediction. A decimal point value,
an interval, a one-sided limit, a text assignment and an unknown are distinct
records. An unknown uncertainty is not zero.

\paragraph{Derivation.}
Let \(s\) denote a complete declared state identity, \(p\) a property and
\(c\) its conditions. An empirical property comparison is admissible only after the
prediction and target have been assigned the same triple:
\begin{equation}
  (s,p,c)_{\mathrm{prediction}}=(s,p,c)_{\mathrm{target}}.
  \label{eq:manual2-property-identity}
\end{equation}
This is a catalog join condition, not a physical observation law. Residuals
between an unselected hypothetical map and a known target remain arithmetic
diagnostics; they do not pass this empirical comparison gate. In particular,
an isotope-unspecified atomic evaluation cannot supply an isotope shift.
A selected line requires its endpoints; a field-dependent shift also requires
the field, state assignment and applicable regime. No magnetic sublevel is
created from an unresolved total angular momentum label.

The constructed histories of \secref{manual2:all-element-campaign} retain their
formation witnesses and relational observables. Assigning an atomic state to a
history and assigning an emitted phase or an energy to that state are additional
maps. Distinct positive maps satisfying the declared native and dimensional
constraints can assign different values to the same history. These typed
witnesses show that those constraints alone do not select a physical value. A common calibration
anchor supplies a scale only after the relevant dimensionless relation has
been established.

\paragraph{Specialization.}
The pilot comprises hydrogen-1, hydrogen-2, hydrogen-3, helium-3, helium-4 and
caesium-133. Atomic identities with unspecified isotope remain separate from these
pilot assignments. The element extension covers atomic numbers 1--118. Its 500 entities carry
4\,354 property records: 1\,067 numeric values, 267 source text assignments,
2\,756 unknowns and 264 inapplicable entries. Broad electronic configuration,
ionization, affinity, level and line tables remain unacquired; their unknown
entries do not assert that physical measurements are unavailable. Neutral-atom masses,
decay channels and state assignments apply to the curated isotope set of
\secref{manual2:isotope-survival}, extended by the pilot identities.

Standard atomic weight and isotope abundance describe a stated material or
terrestrial composition convention. They remain separate from the mass of an
individual isotope. Mass excess, binding energy, separation energy and decay
energy retain their source definitions and electron conventions. The separate
24-row mass calculation retains 15 numerical results and nine missing-input
entries. For binding energy \(B=A b\), the exact mass number scales both the
reported value per nucleon and its quoted uncertainty magnitude; it does not
create \(A\) independent measurements.

For a mass combination \(Q=\sum_i c_iX_i\), repeated occurrences of the same
source quantity are combined first. A compatible covariance matrix gives
\(\operatorname{Var}(Q)=c^{\mathsf T}\Sigma c\), the affine covariance rule of
\mainsection{prop:m3-finite-law}. AME2020 identifies its quoted uncertainties
as one standard deviation\cite[p.~030002-2]{iupac-ame-i}. If their covariance
is unspecified, write these magnitudes as \(u_i\). With
\(w_i=|c_i|u_i\) and \(S=\sum_iw_i\), the possible standard deviation lies
between \(\max(0,2\max_iw_i-S)\) and \(S\). This conditional range is neither
a hard error interval nor an adopted covariance or complete physical budget.

Deuterium illustrates the shared inputs:
\(S_n(\mathrm D)=S_p(\mathrm D)
=\mathrm{ME}(\mathrm H)+\mathrm{ME}(n)-\mathrm{ME}(\mathrm D)
=2224.566229\,\mathrm{keV}\).
The two expressions are the same source-derived quantity, not independent
observations. Their conditional standard-deviation range is
\([0.000412,0.000468]\,\mathrm{keV}\). The separately rounded binding-energy
route gives \(2224.5662\,\mathrm{keV}\); both source reductions are retained.

A separate calculation combines the precise atomic masses of AME2020
Table~A\cite[Table~A]{iupac-ame} with its published covariance. Six light-nuclide
reductions retain their full mutual covariance in \((10^{-9}\,\mathrm u)^2\).
The chosen compatibility criterion requires the covariance's standard deviation
to lie within half a last displayed digit of the quoted marginal uncertainty.
Caesium-133 fails this criterion, leaving its three reductions unresolved.
The uncertain mass-to-energy coefficient is not applied; its dependence on
the mass inputs remains unassigned.

For the retained pilot histories, write \(q_i\) for each generic bounded
cadence in its configuration-specific window. One deliberately hypothetical
cross-configuration bridge is
\begin{equation}
 \widetilde\nu_i=\nu_{\mathrm{Cs}}\frac{q_i}{q_{\mathrm{Cs}}}
 \left(1+\lambda\frac{55-Z_i}{54}\right),\qquad \lambda\geq0.
 \label{eq:manual2-property-bridge-family}
\end{equation}
Here \(\nu_{\mathrm{Cs}}\) is the retained defining frequency. This choice
normalizes Cs for every \(\lambda\) but does not identify the generic cadence
with an atomic transition. The executed hydrogen-isotope cadences are
\(q_{\mathrm H}=1/3\), \(q_{\mathrm D}=1/22\) and
\(q_{\mathrm T}=1/35\). Since all three have \(Z=1\), their shared factor
cancels in D/H and T/H, fixing those ratios throughout this stipulated family.

\paragraph{Calibration.}
The existing hydrogen and caesium records retain their exposure and roles from
\secref{manual2:hyperfine-application}. A defining frequency is exact as a
convention; realization uncertainty belongs to the realization. Already seen
spectroscopic targets remain retrospective. A frozen model protocol fixes its
inputs, common calibration and comparison domain before the separate target
join. Missing state or observation maps leave the corresponding prediction
undefined; they do not produce a zero residual.

\paragraph{Falsification.}
A charge mismatch, substituted isotope, changed transition endpoint or removed
source qualification invalidates a comparison. A new target value must not
alter the generated native history or prediction inputs. Matched feedback and
passive-observation controls distinguish a dynamical intervention from a
display operation. The shared finite dynamics and its state correspondence
must support every claimed observable without an independently tuned gain for
each property. In the stipulated pilot bridge family, the common gain cancels
between hydrogen isotopes. The resulting D/H and T/H ratios are respectively
\(3/22\) and \(3/35\), outside the declared retrospective envelopes formed
from the retained quoted uncertainty magnitudes. This rejects that bridge family
under its stated diagnostic criterion. Those envelopes are not a joint
confidence region; the result supplies neither a completed atomic observation
law nor a rejection of other AOD models.

\aodprovenance{The additive catalog and coverage matrix reside in
\path{manual2/data/m1/properties/}. Its model protocol and separately
generated comparisons retain the distinction between a native construction,
an application choice and an empirical observable. The original element ledger
and its 126 reference positions remain those of
\appref{manual2:app:element-ledger}.}

\aodliteraturenote{The offline atomic adoption interface addresses the NIST Atomic Spectra
Database\cite{nist-asd}. Acquired neutral-atom masses use AME2020\cite{iupac-ame};
spin, abundance and decay assignments use NUBASE2020\cite{isotopes-nubase2020},
with dated moment and radius records from ENSDF\cite{isotopes-ensdf}. The retained source manifests specify the acquired
files, locators and qualifications for each adopted value.}
